% chapters/ch07.tex
% 第~
\ref{chap:7}~章 \index{HilleYosidadingli@Hille--Yosida 定理 (Hille--Yosida theorem)}
Hille--Yosida 定理
% 源: 原书 PDF 第 189--207 页 (印刷页 181--199)

\chapter{Hille--Yosida 定理}
\label{chap:7}

\section{\index{jidadandiaosuanzi@极大单调算子 (maximal monotone operator)}
极大单调算子的定义与初等性质}
\label{sec:7.1}

本章自始至终用 $H$ 表示 Hilbert 空间.

\begin{definition}
称无界线性算子 $A:D(A)\subset H\to H$ 是\emph{单调的},\footnote{有些作者称
$A$ 是\emph{增生算子} (accretive), 或称 $-A$ 是\emph{耗散算子}
(dissipative).} 若它满足
\begin{equation*}
(Av,v)\ge0\quad\forall v\in D(A).
\end{equation*}
若此外还有 $R(I+A)=H$, 即
\begin{equation*}
\forall f\in H\ \ \exists u\in D(A)\ \text{使得}\ u+Au=f,
\end{equation*}
则称 $A$ 是\emph{极大单调的}.
\end{definition}

\begin{proposition}
\label{prop:7.1}
设 $A$ 为极大单调算子.则
\begin{enumerate}[label=(\alph*), leftmargin=2.2em, itemsep=0.1em]
	\item $D(A)$ 在 $H$ 中稠密,
	\item $A$ 是闭算子,
	\item 对每个 $\lambda>0$, $(I+\lambda A)$ 是由 $D(A)$ 到 $H$ 的双射,
	      $(I+\lambda A)^{-1}$ 是有界算子, 且
	      $\norm{(I+\lambda A)^{-1}}_{\LL(H)}\le1$.
\end{enumerate}
\end{proposition}

\begin{proof}
\begin{enumerate}[label=(\alph*), leftmargin=2.2em, itemsep=0.3em]
	\item 设 $f\in H$ 满足 $(f,v)=0\ \forall v\in D(A)$.我们断言 $f=0$.
	      事实上, 存在某个 $v_0\in D(A)$ 使得 $v_0+Av_0=f$.我们有
	      \begin{equation*}
	      0=(f,v_0)=\abs{v_0}^2+(Av_0,v_0)\ge\abs{v_0}^2.
	      \end{equation*}
	      于是 $v_0=0$, 从而 $f=0$.
	\item 首先注意: 给定任何 $f\in H$, 存在唯一的 $u\in D(A)$ 使得
	      $u+Au=f$, 因为若 $u'$ 是另一个解, 则
	      \begin{equation*}
	      u-u'+A(u-u')=0.
	      \end{equation*}
	      与 $(u-u')$ 作标量积并利用单调性, 可见 $u-u'=0$.其次注意
	      $\abs{u}\le\abs{f}$, 因为
	      $\abs{u}^2+(Au,u)=(f,u)\ge\abs{u}^2$.因此映射 $f\mapsto u$
	      (记作 $(I+A)^{-1}$) 是由 $H$ 到自身的有界线性算子, 且
	      $\norm{(I+A)^{-1}}_{\LL(H)}\le1$.下面证明 $A$ 是闭算子.设
	      $(u_n)$ 为 $D(A)$ 中的序列, 使得 $u_n\to u$ 且 $Au_n\to f$.
	      我们要验证 $u\in D(A)$ 且 $Au=f$.但 $u_n+Au_n\to u+f$, 因此
	      \begin{equation*}
	      u_n=(I+A)^{-1}(u_n+Au_n)\to(I+A)^{-1}(u+f).
	      \end{equation*}
	      于是 $u=(I+A)^{-1}(u+f)$, 即 $u\in D(A)$ 且 $u+Au=u+f$.
	\item 我们将证明: 若对某个 $\lambda_0>0$ 有 $R(I+\lambda_0A)=H$, 则对每个
	      $\lambda>\lambda_0/2$ 都有 $R(I+\lambda A)=H$.首先注意——如 (b) 中
	      那样——对每个 $f\in H$ 存在唯一的 $u\in D(A)$ 使得
	      $u+\lambda_0Au=f$.此外, 映射 $f\mapsto u$ (记作
	      $(I+\lambda_0A)^{-1}$) 是有界线性算子, 且
	      $\norm{(I+\lambda_0A)^{-1}}_{\LL(H)}\le1$.我们试图求解方程
	      \begin{equation}
	      u+\lambda Au=f\quad\text{其中}\ \lambda>0.
	      \label{eq:7.1}
	      \end{equation}
	      方程 \eqref{eq:7.1} 可写成
	      \begin{equation*}
	      u+\lambda_0Au=\frac{\lambda_0}{\lambda}f
	      +\left(1-\frac{\lambda_0}{\lambda}\right)u,
	      \end{equation*}
	      或者等价地
	      \begin{equation}
	      u=(I+\lambda_0A)^{-1}
	      \left[\frac{\lambda_0}{\lambda}f
	      +\left(1-\frac{\lambda_0}{\lambda}\right)u\right].
	      \label{eq:7.2}
	      \end{equation}
	      若 $\left|1-\frac{\lambda_0}{\lambda}\right|<1$, 即
	      $\lambda>\lambda_0/2$, 则可应用压缩映射原理 (定理  \ref{thm:5.7}),
	      推出 \eqref{eq:7.2} 有解.

	      结论 (c) 容易由归纳法得出: 由于 $I+A$ 满射, 故对每个
	      $\lambda>1/2$, $I+\lambda A$ 满射, 从而对每个 $\lambda>1/4$
	      也满射, 依此类推.
\end{enumerate}
\end{proof}

\begin{remark}
\label{rem:7.1}
若 $A$ 极大单调, 则对每个 $\lambda>0$, $\lambda A$ 也极大单调.然而, 若
$A$ 与 $B$ 都是极大单调算子, 则定义在 $D(A)\cap D(B)$ 上的 $A+B$ 不必是
极大单调的.
\end{remark}

\begin{definition}
设 $A$ 为极大单调算子.对每个 $\lambda>0$, 令
\begin{equation*}
J_\lambda=(I+\lambda A)^{-1}
\quad\text{与}\quad
A_\lambda=\frac1\lambda(I-J_\lambda);
\end{equation*}
$J_\lambda$ 称为 $A$ 的\emph{预解式}, $A_\lambda$ 称为 $A$ 的
\emph{\index{Yosidabijin@Yosida 逼近 (Yosida approximation)}
Yosida 逼近} (或\emph{正则化}).记住
$\norm{J_\lambda}_{\LL(H)}\le1$.
\end{definition}

\begin{proposition}
\label{prop:7.2}
设 $A$ 为极大单调算子.则
\begin{align*}
A_\lambda v&=A(J_\lambda v)
&&\forall v\in H\ \text{及}\ \forall\lambda>0,
\tag*{(a$_1$)}\\
A_\lambda v&=J_\lambda(Av)
&&\forall v\in D(A)\ \text{及}\ \forall\lambda>0,
\tag*{(a$_2$)}\\
\abs{A_\lambda v}&\le\abs{Av}
&&\forall v\in D(A)\ \text{及}\ \forall\lambda>0,
\tag*{(b)}\\
\lim_{\lambda\to0}J_\lambda v&=v
&&\forall v\in H,
\tag*{(c)}\\
\lim_{\lambda\to0}A_\lambda v&=Av
&&\forall v\in D(A),
\tag*{(d)}\\
(A_\lambda v,v)&\ge0
&&\forall v\in H\ \text{及}\ \forall\lambda>0,
\tag*{(e)}\\
\abs{A_\lambda v}&\le(1/\lambda)\abs{v}
&&\forall v\in H\ \text{及}\ \forall\lambda>0.
\tag*{(f)}
\end{align*}
\end{proposition}

\begin{proof}
(a$_1$) 可写成 $v=(J_\lambda v)+\lambda A(J_\lambda v)$, 这正是 $J_\lambda v$
的定义.

(a$_2$) 由 (a$_1$) 我们有
\begin{equation*}
A_\lambda v+A(v-J_\lambda v)=Av,
\end{equation*}
即
\begin{equation*}
A_\lambda v+\lambda A(A_\lambda v)=Av,
\end{equation*}
这意味着 $A_\lambda v=(I+\lambda A)^{-1}Av$.

(b) 容易由 (a$_2$) 得出.

(c) 先设 $v\in D(A)$.则
\begin{equation*}
\abs{v-J_\lambda v}=\lambda\abs{A_\lambda v}\le\lambda\abs{Av}
\quad\text{(由 (b))},
\end{equation*}
从而 $\lim_{\lambda\to0}J_\lambda v=v$.

现在设 $v$ 是 $H$ 中的一般元素.给定任何 $\varepsilon>0$, 存在某个
$v_1\in D(A)$ 使得 $\abs{v-v_1}\le\varepsilon$ (因为由命题  \ref{prop:7.1},
$D(A)$ 在 $H$ 中稠密).我们有
\begin{equation*}
\abs{J_\lambda v-v}\le\abs{J_\lambda v-J_\lambda v_1}
+\abs{J_\lambda v_1-v_1}+\abs{v_1-v}
\le2\abs{v-v_1}+\abs{J_\lambda v_1-v_1}
\le2\varepsilon+\abs{J_\lambda v_1-v_1}.
\end{equation*}
于是
\begin{equation*}
\limsup_{\lambda\to0}\abs{J_\lambda v-v}\le2\varepsilon
\quad\forall\varepsilon>0,
\end{equation*}
从而
\begin{equation*}
\lim_{\lambda\to0}\abs{J_\lambda v-v}=0.
\end{equation*}

(d) 这是 (a$_2$) 与 (c) 的推论.

(e) 我们有
\begin{equation*}
(A_\lambda v,v)=(A_\lambda v,v-J_\lambda v)+(A_\lambda v,J_\lambda v)
=\lambda\abs{A_\lambda v}^2+(A(J_\lambda v),J_\lambda v),
\end{equation*}
从而
\begin{equation}
(A_\lambda v,v)\ge\lambda\abs{A_\lambda v}^2.
\label{eq:7.3}
\end{equation}

(f) 这是 \eqref{eq:7.3} 与 Cauchy--Schwarz 不等式的推论.
\end{proof}

\begin{remark}
\label{rem:7.2}
命题  \ref{prop:7.2} 说明 $(A_\lambda)_{\lambda>0}$ 是一族有界算子, 当
$\lambda\to0$ 时它们「逼近」无界算子 $A$.这种逼近方式将被非常频繁地使用.
当然, 一般说来 $\norm{A_\lambda}_{\LL(H)}$ 当 $\lambda\to0$ 时会「爆炸」.
\end{remark}

\section{发展方程 \texorpdfstring{$du/dt+Au=0$}{du/dt+Au=0} 在
\texorpdfstring{$[0,+\infty)$}{[0,+∞)} 上的解.
\texorpdfstring{$u(0)=u_0$}{u(0)=u0}. 存在性与唯一性}
\label{sec:7.2}

我们先从一个非常经典的结果开始:

\begin{keytheorem}[Cauchy, Lipschitz, Picard]
\label{thm:7.3}
设 $E$ 为 Banach 空间, $F:E\to E$ 为 Lipschitz 映射, 即存在常数 $L$,
使得
\begin{equation*}
\norm{Fu-Fv}\le L\norm{u-v}\quad\forall u,v\in E.
\end{equation*}
则给定任何 $u_0\in E$, 存在唯一的解 $u\in C^1([0,+\infty);E)$ 满足
\begin{equation}
\begin{cases}
\dfrac{du}{dt}(t)=Fu(t) & \text{在}\ [0,+\infty)\ \text{上},\\
u(0)=u_0.
\end{cases}
\label{eq:7.4}
\end{equation}
$u_0$ 称为\emph{初值}.
\end{keytheorem}

\begin{proof}
\emph{存在性.} 求解 \eqref{eq:7.4} 等价于求某个
$u\in C([0,+\infty);E)$ 满足积分方程
\begin{equation}
u(t)=u_0+\int_0^tF(u(s))\,ds.
\label{eq:7.5}
\end{equation}
给定 $k>0$ (稍后确定), 令
\begin{equation*}
X=\Bigl\{u\in C([0,+\infty);E);\ \sup_{t\ge0}e^{-kt}\norm{u(t)}<\infty\Bigr\}.
\end{equation*}
容易验证 $X$ 关于范数
\begin{equation*}
\norm{u}_X=\sup_{t\ge0}e^{-kt}\norm{u(t)}
\end{equation*}
是 Banach 空间.对每个 $u\in X$, 由
\begin{equation*}
(\mathcal{T}u)(t)=u_0+\int_0^tF(u(s))\,ds
\end{equation*}
定义的函数 $\mathcal{T}u$ 也属于 $X$.此外, 我们有
\begin{equation*}
\norm{\mathcal{T}u-\mathcal{T}v}_X\le\frac Lk\norm{u-v}_X
\quad\forall u,v\in X.
\end{equation*}
取 $k>L$, 可见 $\mathcal{T}$ 在 $X$ 中有 (唯一) 不动点 $u$, 它是
\eqref{eq:7.5} 的解.

\emph{唯一性.} 设 $u$ 与 $u'$ 是 \eqref{eq:7.4} 的两个解, 令
\begin{equation*}
\phi(t)=\norm{u(t)-u'(t)}.
\end{equation*}
由 \eqref{eq:7.5} 推出
\begin{equation*}
\phi(t)\le L\int_0^t\phi(s)\,ds\quad\forall t\ge0,
\end{equation*}
从而 $\phi\equiv0$.
\end{proof}

前一个定理在研究常微分方程时极其有用.然而在研究偏微分方程时它用处不大.
我们下一个结果在求解发展型偏微分方程时是非常强有力的工具; 见第~
\ref{chap:10}~章.

\begin{keytheorem}[Hille--Yosida]
\label{thm:7.4}
设 $A$ 为极大单调算子.则给定任何 $u_0\in D(A)$, 存在唯一的函数\footnote{空间
$D(A)$ 赋以图范数 $\abs{v}+\abs{Av}$, 或与之等价的 Hilbert 范数
$(\abs{v}^2+\abs{Av}^2)^{1/2}$.}
\begin{equation*}
u\in C^1([0,+\infty);H)\cap C([0,+\infty);D(A))
\end{equation*}
满足
\begin{equation}
\begin{cases}
\dfrac{du}{dt}+Au=0 & \text{在}\ [0,+\infty)\ \text{上},\\
u(0)=u_0.
\end{cases}
\label{eq:7.6}
\end{equation}
此外,
\begin{equation*}
\abs{u(t)}\le\abs{u_0}
\quad\text{且}\quad
\left|\frac{du}{dt}(t)\right|=\abs{Au(t)}\le\abs{Au_0}
\quad\forall t\ge0.
\end{equation*}
\end{keytheorem}

\begin{remark}
\label{rem:7.3}
定理  \ref{thm:7.4} 的主要意义在于, 它把「发展问题」的研究归结为「定态方程」
$u+Au=f$ 的研究 (假定我们已经知道 $A$ 单调, 而这在实践中容易验证).
\end{remark}

\begin{proof}
证明分为六步.

\emph{第 1 步: 唯一性.} 设 $u$ 与 $u'$ 是 \eqref{eq:7.6} 的两个解.我们有
\begin{equation*}
\left(\frac{d}{dt}(u-u'),(u-u')\right)
=-(A(u-u'),u-u')\le0.
\end{equation*}
但\footnote{记住: 若 $\phi\in C^1([0,+\infty);H)$, 则
$\abs{\phi}^2\in C^1([0,+\infty);\R)$, 且
$\dfrac{d}{dt}\abs{\phi}^2=2\left(\dfrac{d\phi}{dt},\phi\right)$.}
\begin{equation*}
\frac12\frac{d}{dt}\abs{u(t)-u'(t)}^2
=\left(\frac{d}{dt}(u(t)-u'(t)),u(t)-u'(t)\right).
\end{equation*}
因此函数 $t\mapsto\abs{u(t)-u'(t)}$ 在 $[0,+\infty)$ 上不增.由于
$\abs{u(0)-u'(0)}=0$, 推出
\begin{equation*}
\abs{u(t)-u'(t)}=0\quad\forall t\ge0.
\end{equation*}

证明存在性的主要想法是: 在 \eqref{eq:7.6} 中用 $A_\lambda$ 代替 $A$, 对
近似问题应用定理  \ref{thm:7.3}, 然后利用各种与 $\lambda$ 无关的估计令
$\lambda\to0$ 取极限.于是设 $u_\lambda$ 为问题
\begin{equation}
\begin{cases}
\dfrac{du_\lambda}{dt}+A_\lambda u_\lambda=0
& \text{在}\ [0,+\infty)\ \text{上},\\
u_\lambda(0)=u_0\in D(A)
\end{cases}
\label{eq:7.7}
\end{equation}
的解.

\emph{第 2 步:} 我们有估计
\begin{equation}
\abs{u_\lambda(t)}\le\abs{u_0}
\qquad\forall t\ge0,\ \forall\lambda>0,
\label{eq:7.8}
\end{equation}
\begin{equation}
\left|\frac{du_\lambda}{dt}(t)\right|=\abs{A_\lambda u_\lambda(t)}
\le\abs{Au_0}
\qquad\forall t\ge0,\ \forall\lambda>0.
\label{eq:7.9}
\end{equation}
它们直接由下述引理以及 $\abs{A_\lambda u_0}\le\abs{Au_0}$ 得出.

\begin{lemma}
\label{lem:7.1}
设 $w\in C^1([0,+\infty);H)$ 满足
\begin{equation}
\frac{dw}{dt}+A_\lambda w=0\quad\text{在}\ [0,+\infty)\ \text{上}.
\label{eq:7.10}
\end{equation}
则函数 $t\mapsto\abs{w(t)}$ 与 $t\mapsto
\left|\dfrac{dw}{dt}(t)\right|=\abs{A_\lambda w(t)}$ 在 $[0,+\infty)$ 上不增.
\end{lemma}

\begin{proof}
我们有
\begin{equation*}
\left(\frac{dw}{dt},w\right)+(A_\lambda w,w)=0.
\end{equation*}
由命题  \ref{prop:7.2}(e) 我们知道 $(A_\lambda w,w)\ge0$, 从而
$\frac12\frac{d}{dt}\abs{w}^2\le0$, 故 $\abs{w(t)}$ 不增.另一方面, 由于
$A_\lambda$ 是线性有界算子, 由 \eqref{eq:7.10} 可 (用归纳法) 推出
$w\in C^\infty([0,+\infty);H)$, 并且
\begin{equation*}
\frac{d}{dt}\left(\frac{dw}{dt}\right)
+A_\lambda\left(\frac{dw}{dt}\right)=0.
\end{equation*}
把上述事实应用于 $\dfrac{dw}{dt}$, 可见
$\left|\dfrac{dw}{dt}(t)\right|$ 不增.事实上, 对任意阶 $k$, 函数
$\left|\dfrac{d^kw}{dt^k}(t)\right|$ 都不增.
\end{proof}

\emph{第 3 步:} 这里我们将证明: 对每个 $t\ge0$, 当 $\lambda\to0$ 时
$u_\lambda(t)$ 收敛于某个极限, 记作 $u(t)$.此外, 这个收敛在每个有界区间
$[0,T]$ 上是一致的.

对每个 $\lambda,\mu>0$, 我们有
\begin{equation*}
\frac{du_\lambda}{dt}-\frac{du_\mu}{dt}
+A_\lambda u_\lambda-A_\mu u_\mu=0,
\end{equation*}
从而
\begin{equation}
\frac12\frac{d}{dt}\abs{u_\lambda(t)-u_\mu(t)}^2
+(A_\lambda u_\lambda(t)-A_\mu u_\mu(t),u_\lambda(t)-u_\mu(t))=0.
\label{eq:7.11}
\end{equation}
为简单起见略去 $t$, 我们写出
\begin{align}
&(A_\lambda u_\lambda-A_\mu u_\mu,u_\lambda-u_\mu)\notag\\
&\quad=(A_\lambda u_\lambda-A_\mu u_\mu,
u_\lambda-J_\lambda u_\lambda+J_\lambda u_\lambda
-J_\mu u_\mu+J_\mu u_\mu-u_\mu)\notag\\
&\quad=(A_\lambda u_\lambda-A_\mu u_\mu,\lambda A_\lambda u_\lambda
-\mu A_\mu u_\mu)+(A(J_\lambda u_\lambda-J_\mu u_\mu),
J_\lambda u_\lambda-J_\mu u_\mu)\notag\\
&\quad\ge(A_\lambda u_\lambda-A_\mu u_\mu,\lambda A_\lambda u_\lambda
-\mu A_\mu u_\mu).
\label{eq:7.12}
\end{align}
由 \eqref{eq:7.9}, \eqref{eq:7.11} 与 \eqref{eq:7.12} 推出
\begin{equation*}
\frac12\frac{d}{dt}\abs{u_\lambda-u_\mu}^2
\le2(\lambda+\mu)\abs{Au_0}^2.
\end{equation*}
对此不等式积分, 得到
\begin{equation*}
\abs{u_\lambda(t)-u_\mu(t)}^2\le4(\lambda+\mu)t\abs{Au_0}^2,
\end{equation*}
即
\begin{equation}
\abs{u_\lambda(t)-u_\mu(t)}\le2\sqrt{(\lambda+\mu)t}\,\abs{Au_0}.
\label{eq:7.13}
\end{equation}
由此推出, 对每个固定的 $t\ge0$, 当 $\lambda\to0$ 时 $u_\lambda(t)$ 是
Cauchy 列, 从而收敛于某个极限, 记作 $u(t)$.在 \eqref{eq:7.13} 中令
$\mu\to0$ 取极限, 我们有
\begin{equation*}
\abs{u_\lambda(t)-u(t)}\le2\sqrt{\lambda t}\,\abs{Au_0}.
\end{equation*}
因此, 收敛在 $t$ 上于每个有界区间 $[0,T]$ 一致, 故
$u\in C([0,+\infty);H)$.

\emph{第 4 步:} 这里再假设 $u_0\in D(A^2)$, 即 $u_0\in D(A)$ 且
$Au_0\in D(A)$, 我们证明当 $\lambda\to0$ 时
$\dfrac{du_\lambda}{dt}(t)$ 收敛于某个极限, 且收敛在每个有界区间 $[0,T]$ 上
一致.

令 $v_\lambda=\dfrac{du_\lambda}{dt}$, 于是
$\dfrac{dv_\lambda}{dt}+A_\lambda v_\lambda=0$.沿用第 3 步的论证, 可见
\begin{equation}
\frac12\frac{d}{dt}\abs{v_\lambda-v_\mu}^2
\le(\abs{A_\lambda v_\lambda}+\abs{A_\mu v_\mu})
(\lambda\abs{A_\lambda v_\lambda}+\mu\abs{A_\mu v_\mu}).
\label{eq:7.14}
\end{equation}
由引理  \ref{lem:7.1}, 我们有
\begin{equation}
\abs{A_\lambda v_\lambda(t)}\le\abs{A_\lambda v_\lambda(0)}
=\abs{A_\lambda A_\lambda u_0},
\label{eq:7.15}
\end{equation}
类似地
\begin{equation}
\abs{A_\mu v_\mu(t)}\le\abs{A_\mu v_\mu(0)}
=\abs{A_\mu A_\mu u_0}.
\label{eq:7.16}
\end{equation}
最后, 由于 $Au_0\in D(A)$, 我们得到
\begin{equation*}
A_\lambda A_\lambda u_0=J_\lambda AJ_\lambda Au_0
=J_\lambda J_\lambda AAu_0=J_\lambda^2A^2u_0,
\end{equation*}
从而
\begin{equation}
\abs{A_\lambda A_\lambda u_0}\le\abs{A^2u_0},
\qquad
\abs{A_\mu A_\mu u_0}\le\abs{A^2u_0}.
\label{eq:7.17}
\end{equation}
合并 \eqref{eq:7.14}, \eqref{eq:7.15}, \eqref{eq:7.16} 与
\eqref{eq:7.17}, 我们得到
\begin{equation*}
\frac12\frac{d}{dt}\abs{v_\lambda-v_\mu}^2
\le2(\lambda+\mu)\abs{A^2u_0}^2.
\end{equation*}
如第 3 步那样, 我们得出结论: 当 $\lambda\to0$ 时
$v_\lambda(t)=\dfrac{du_\lambda}{dt}(t)$ 收敛于某个极限, 且收敛在每个有界
区间 $[0,T]$ 上一致.

\emph{第 5 步:} 假设 $u_0\in D(A^2)$, 这里证明 $u$ 是 \eqref{eq:7.6} 的解.

由第 3 步与第 4 步, 我们知道对一切 $T<\infty$,
\begin{equation*}
\begin{cases}
u_\lambda(t)\to u(t), & \text{当}\ \lambda\to0,\ \text{在}\ [0,T]\ \text{上一致},\\
\dfrac{du_\lambda}{dt}(t)\ \text{当}\ \lambda\to0\ \text{时收敛},
& \text{在}\ [0,T]\ \text{上一致}.
\end{cases}
\end{equation*}
由此容易推出 $u\in C^1([0,+\infty);H)$, 且当 $\lambda\to0$ 时
$\dfrac{du_\lambda}{dt}(t)\to\dfrac{du}{dt}(t)$ 在 $[0,T]$ 上一致.把
\eqref{eq:7.7} 改写为
\begin{equation}
\frac{du_\lambda}{dt}(t)+A(J_\lambda u_\lambda(t))=0.
\label{eq:7.18}
\end{equation}
注意当 $\lambda\to0$ 时 $J_\lambda u_\lambda(t)\to u(t)$, 因为
\begin{equation*}
\abs{J_\lambda u_\lambda(t)-u(t)}
\le\abs{J_\lambda u_\lambda(t)-J_\lambda u(t)}
+\abs{J_\lambda u(t)-u(t)}
\le\abs{u_\lambda(t)-u(t)}+\abs{J_\lambda u(t)-u(t)}\to0.
\end{equation*}
利用 $A$ 有闭图像这一事实, 由 \eqref{eq:7.18} 推出
$u(t)\in D(A)\ \forall t\ge0$, 且
\begin{equation*}
\frac{du}{dt}(t)+Au(t)=0.
\end{equation*}
最后, 由于 $u\in C^1([0,+\infty);H)$, 函数 $t\mapsto Au(t)$ 由
$[0,+\infty)$ 到 $H$ 连续, 故 $u\in C([0,+\infty);D(A))$.于是我们得到了
\eqref{eq:7.6} 的一个解, 且此外还满足
\begin{equation*}
\abs{u(t)}\le\abs{u_0}\quad\forall t\ge0
\quad\text{与}\quad
\left|\frac{du}{dt}(t)\right|=\abs{Au(t)}\le\abs{Au_0}
\quad\forall t\ge0.
\end{equation*}

\emph{第 6 步:} 这里我们完成定理的证明.

我们将使用下述引理.

\begin{lemma}
\label{lem:7.2}
设 $u_0\in D(A)$.则 $\forall\varepsilon>0$, $\exists u'_0\in D(A^2)$,
使得 $\abs{u_0-u'_0}<\varepsilon$ 且 $\abs{Au_0-Au'_0}<\varepsilon$.换言之,
$D(A^2)$ 在 $D(A)$ 中稠密 (关于图范数).
\end{lemma}

\begin{proof}[引理 \ref{lem:7.2}]
令 $u'_0=J_\lambda u_0$, 其中 $\lambda>0$ 取得适当小 (稍后确定).我们有
\begin{equation*}
u'_0\in D(A)\quad\text{且}\quad u'_0+\lambda Au'_0=u_0.
\end{equation*}
于是 $Au'_0\in D(A)$, 即 $u'_0\in D(A^2)$.另一方面, 由命题
 \ref{prop:7.2}, 我们知道
\begin{equation*}
\lim_{\lambda\to0}\abs{J_\lambda u_0-u_0}=0,\qquad
\lim_{\lambda\to0}\abs{J_\lambda Au_0-Au_0}=0,
\end{equation*}
且
\begin{equation*}
J_\lambda Au_0=AJ_\lambda u_0.
\end{equation*}
取 $\lambda>0$ 充分小即得所需结论.
\end{proof}

现在我们回到定理  \ref{thm:7.4} 的证明.给定 $u_0\in D(A)$, 我们 (用引理
 \ref{lem:7.2}) 构造 $D(A^2)$ 中的序列 $(u_{0n})$, 使得 $u_{0n}\to u_0$ 且
$Au_{0n}\to Au_0$.由第 5 步, 我们知道问题
\begin{equation}
\begin{cases}
\dfrac{du_n}{dt}+Au_n=0 & \text{在}\ [0,+\infty)\ \text{上},\\
u_n(0)=u_{0n}
\end{cases}
\label{eq:7.19}
\end{equation}
有解 $u_n$.对一切 $t\ge0$, 我们有
\begin{equation*}
\abs{u_n(t)-u_m(t)}\le\abs{u_{0n}-u_{0m}}
\xrightarrow[m,n\to\infty]{}0,
\end{equation*}
\begin{equation*}
\left|\frac{du_n}{dt}(t)-\frac{du_m}{dt}(t)\right|
\le\abs{Au_{0n}-Au_{0m}}
\xrightarrow[m,n\to\infty]{}0.
\end{equation*}
因此
\begin{equation*}
u_n(t)\to u(t)\ \text{在}\ [0,+\infty)\ \text{上一致},
\qquad
\frac{du_n}{dt}(t)\to\frac{du}{dt}(t)
\ \text{在}\ [0,+\infty)\ \text{上一致},
\end{equation*}
其中 $u\in C^1([0,+\infty);H)$.在 \eqref{eq:7.19} 中取极限——利用 $A$ 是
闭算子这一事实——可见 $u(t)\in D(A)$ 且 $u$ 满足 \eqref{eq:7.6}.由
\eqref{eq:7.6} 推出 $u\in C([0,+\infty);D(A))$.
\end{proof}

\begin{remark}
\label{rem:7.4}
设 $u_\lambda$ 为 \eqref{eq:7.7} 的解:
\begin{enumerate}[label=(\alph*), leftmargin=2.2em, itemsep=0.2em]
	\item 假设 $u_0\in D(A)$.我们知道 (由第 3 步), 当 $\lambda\to0$ 时,
	      对每个 $t\ge0$, $u_\lambda(t)$ 收敛于某个极限 $u(t)$.可以直接
	      证明 $u\in C^1([0,+\infty);H)\cap C([0,+\infty);D(A))$, 且它满足
	      \eqref{eq:7.6}.
	\item 只假设 $u_0\in H$.仍可证明当 $\lambda\to0$ 时, 对每个 $t\ge0$,
	      $u_\lambda(t)$ 收敛于某个极限, 记作 $u(t)$.但可能出现这样的情形:
	      这个极限 $u(t)$ 对一切 $t>0$ 都不属于 $D(A)$, 且 $u(t)$ 在
	      $[0,+\infty)$ 上处处不可微.因此 $u(t)$ 不是 \eqref{eq:7.6} 的
	      「古典」解.事实上, 对这样的 $u_0$, 问题 \eqref{eq:7.6} 没有古典
	      解.尽管如此, 我们仍可把 $u(t)$ 看作 \eqref{eq:7.6} 的「广义」解.
	      我们将在第~\ref{sec:7.4}~节看到, 当 $A$ 自伴时这种情况不会发生: 此时对每个
	      $u_0\in H$ (即使 $u_0\notin D(A)$), $u(t)$ 都是 \eqref{eq:7.6}
	      的「古典」解.
\end{enumerate}
\end{remark}

\begin{starremark}[压缩半群]
\label{rem:7.5}
对每个 $t\ge0$, 考虑线性映射 $u_0\in D(A)\mapsto u(t)\in D(A)$, 其中
$u(t)$ 是定理  \ref{thm:7.4} 给出的 \eqref{eq:7.6} 的解.由于
$\abs{u(t)}\le\abs{u_0}$ 且 $D(A)$ 在 $H$ 中稠密, 我们可以按连续性把这个
映射延拓为由 $H$ 到自身的有界算子, 记作 $S_A(t)$.\footnote{或者也可以
用注记 \ref{rem:7.4}(b) 直接在 $H$ 上把 $S_A(t)$ 定义为映射
$u_0\in H\mapsto u(t)\in H$.} 容易验证 $S_A(t)$ 满足下述性质:
\begin{enumerate}[label=(\alph*), leftmargin=2.2em, itemsep=0.2em]
	\item 对每个 $t\ge0$, $S_A(t)\in\LL(H)$ 且
	      $\norm{S_A(t)}_{\LL(H)}\le1$,
	\item
	      \begin{equation*}
	      \begin{cases}
	      S_A(t_1+t_2)=S_A(t_1)\circ S_A(t_2) & \forall t_1,t_2\ge0,\\
	      S_A(0)=I,
	      \end{cases}
	      \end{equation*}
	\item
	      \begin{equation*}
	      \lim_{\substack{t\to0\\ t>0}}\abs{S_A(t)u_0-u_0}=0
	      \quad\forall u_0\in H.
	      \end{equation*}
\end{enumerate}
这样一族依赖于参数 $t\ge0$, 由 $H$ 到自身的算子 $\{S(t)\}_{t\ge0}$ 若满足
(a), (b), (c), 则称为\emph{连续压缩半群}.

Hille 与 Yosida 的一个值得注意的结果断言: 反之, 给定 $H$ 上的连续压缩半群
$S(t)$, 存在唯一的极大单调算子 $A$, 使得
$S(t)=S_A(t)\ \forall t\ge0$.这就在极大单调算子与连续压缩半群之间建立了
一一对应.(证明见第~
\ref{chap:7}~章评注中引用的文献.)
\end{starremark}

\begin{keyremark}
\label{rem:7.6}
设 $A$ 为极大单调算子, $\lambda\in\R$.问题
\begin{equation*}
\begin{cases}
\dfrac{du}{dt}+Au+\lambda u=0 & \text{在}\ [0,+\infty)\ \text{上},\\
u(0)=u_0
\end{cases}
\end{equation*}
可用下述简单办法归结为问题 \eqref{eq:7.6}.令
\begin{equation*}
v(t)=e^{\lambda t}u(t).
\end{equation*}
则 $v$ 满足
\begin{equation*}
\begin{cases}
\dfrac{dv}{dt}+Av=0 & \text{在}\ [0,+\infty)\ \text{上},\\
v(0)=u_0.
\end{cases}
\end{equation*}
\end{keyremark}

\section{正则性}
\label{sec:7.3}

这里我们将证明: 只要对初值 $u_0$ 作附加假设, 定理  \ref{thm:7.4} 中得到的
\eqref{eq:7.6} 的解 $u$ 就比 $C^1([0,+\infty);H)\cap
C([0,+\infty);D(A))$ 更正则.为此, 我们用归纳法定义空间
\begin{equation*}
D(A^k)=\{v\in D(A^{k-1});\ Av\in D(A^{k-1})\},
\end{equation*}
其中 $k$ 为任何整数, $k\ge2$.容易看出 $D(A^k)$ 关于标量积
\begin{equation*}
(u,v)_{D(A^k)}=\sum_{j=0}^k(A^ju,A^jv)
\end{equation*}
是 Hilbert 空间; 相应的范数为
\begin{equation*}
\abs{u}_{D(A^k)}=\left(\sum_{j=0}^k\abs{A^ju}^2\right)^{1/2}.
\end{equation*}

\begin{theorem}
\label{thm:7.5}
假设对某个整数 $k\ge2$, $u_0\in D(A^k)$.则定理  \ref{thm:7.4} 中得到的
问题 \eqref{eq:7.6} 的解 $u$ 满足
\begin{equation*}
u\in C^{k-j}([0,+\infty);D(A^j))\quad\forall j=0,1,\dots,k.
\end{equation*}
\end{theorem}

\begin{proof}
先设 $k=2$.考虑赋以标量积 $(u,v)_{D(A)}$ 的 Hilbert 空间 $H_1=D(A)$.
容易验证由
\begin{equation*}
\begin{cases}
D(A_1)=D(A^2),\\
A_1u=Au & \text{对}\ u\in D(A_1)
\end{cases}
\end{equation*}
定义的算子 $A_1:D(A_1)\subset H_1\to H_1$ 在 $H_1$ 中极大单调.把定理
 \ref{thm:7.4} 应用于空间 $H_1$ 中的算子 $A_1$, 可见存在函数
\begin{equation*}
u\in C^1([0,+\infty);H_1)\cap C([0,+\infty);D(A_1))
\end{equation*}
使得
\begin{equation*}
\begin{cases}
\dfrac{du}{dt}+A_1u=0 & \text{在}\ [0,+\infty)\ \text{上},\\
u(0)=u_0.
\end{cases}
\end{equation*}
特别地, $u$ 满足 \eqref{eq:7.6}; 由唯一性, 这个 $u$ 就是 \eqref{eq:7.6}
的解.只须再验证 $u\in C^2([0,+\infty);H)$.由于
\begin{equation*}
A\in\LL(H_1,H)\quad\text{且}\quad u\in C([0,+\infty);H_1),
\end{equation*}
推出 $Au\in C^1([0,+\infty);H)$ 且
\begin{equation}
\frac{d}{dt}(Au)=A\left(\frac{du}{dt}\right).
\label{eq:7.20}
\end{equation}
应用 \eqref{eq:7.6}, 可见 $\dfrac{du}{dt}\in C^1([0,+\infty);H)$, 即
$u\in C^2([0,+\infty);H)$, 且
\begin{equation}
\frac{d}{dt}\left(\frac{du}{dt}\right)
+A\left(\frac{du}{dt}\right)=0\quad\text{在}\ [0,+\infty)\ \text{上}.
\label{eq:7.21}
\end{equation}

现在我们转向一般情形 $k\ge3$.我们对 $k$ 归纳: 假设结论对直到 $k-1$ 的
阶数都成立, 并设 $u_0\in D(A^k)$.由前面的分析, 我们知道 \eqref{eq:7.6}
的解 $u$ 属于 $C^2([0,+\infty);H)\cap C^1([0,+\infty);D(A))$, 且 $u$ 满足
\eqref{eq:7.21}.令
\begin{equation*}
v=\frac{du}{dt},
\end{equation*}
我们有
\begin{equation*}
v\in C^1([0,+\infty);H)\cap C([0,+\infty);D(A)),
\end{equation*}
\begin{equation*}
\begin{cases}
\dfrac{dv}{dt}+Av=0 & \text{在}\ [0,+\infty)\ \text{上},\\
v(0)=-Au_0.
\end{cases}
\end{equation*}
换言之, $v$ 是对应于初值 $v_0=-Au_0$ 的 \eqref{eq:7.6} 的解.由于
$v_0\in D(A^{k-1})$, 由归纳假设我们知道
\begin{equation}
v\in C^{k-1-j}([0,+\infty);D(A^j))\quad\forall j=0,1,\dots,k-1,
\label{eq:7.22}
\end{equation}
即
\begin{equation*}
u\in C^{k-j}([0,+\infty);D(A^j))\quad\forall j=0,1,\dots,k-1.
\end{equation*}
只须再验证
\begin{equation}
u\in C([0,+\infty);D(A^k)).
\label{eq:7.23}
\end{equation}
在 \eqref{eq:7.22} 中取 $j=k-1$, 可见
\begin{equation}
\frac{du}{dt}\in C([0,+\infty);D(A^{k-1})).
\label{eq:7.24}
\end{equation}
由 \eqref{eq:7.24} 与方程 \eqref{eq:7.6} 推出
\begin{equation*}
Au\in C([0,+\infty);D(A^{k-1})),
\end{equation*}
即 \eqref{eq:7.23}.
\end{proof}

\section{自伴情形}
\label{sec:7.4}

设 $A:D(A)\subset H\to H$ 为无界线性算子, $\overline{D(A)}=H$.把
$H^\star$ 与 $H$ 等同, 我们可以把 $A^\star$ 看作 $H$ 中的无界线性算子.

\begin{definition}
称
\begin{itemize}[leftmargin=1.6em, itemsep=0.15em]
	\item $A$ 是\emph{对称的}, 若 $(Au,v)=(u,Av)\ \forall u,v\in D(A)$,
	\item $A$ 是\emph{自伴的}, 若 $D(A^\star)=D(A)$ 且 $A^\star=A$.
\end{itemize}
\end{definition}

\begin{remark}
\label{rem:7.7}
对有界算子, 对称算子与自伴算子的概念一致.然而, 若 $A$ 无界, 则对称算子与
自伴算子之间有微妙的差别.显然, 任何自伴算子都是对称的.逆命题不成立: 算子
$A$ 对称当且仅当 $A\subset A^\star$, 即 $D(A)\subset D(A^\star)$ 且在
$D(A)$ 上 $A^\star=A$.可能出现 $A$ 对称而 $D(A)\ne D(A^\star)$ 的情形.
我们下一个结果表明: 若 $A$ 极大单调, 则
\begin{equation*}
(A\ \text{对称})\iff(A\ \text{自伴}).
\end{equation*}
\end{remark}

\begin{proposition}
\label{prop:7.6}
设 $A$ 为极大单调的对称算子.则 $A$ 自伴.
\end{proposition}

\begin{proof}
令 $J_1=(I+A)^{-1}$.我们先证明 $J_1$ 自伴.由于 $J_1\in\LL(H)$, 只须
验证
\begin{equation}
(J_1u,v)=(u,J_1v)\quad\forall u,v\in H.
\label{eq:7.25}
\end{equation}
令 $u_1=J_1u$, $v_1=J_1v$, 于是
\begin{equation*}
u_1+Au_1=u,\qquad v_1+Av_1=v.
\end{equation*}
由于按假设 $(u_1,Av_1)=(Au_1,v_1)$, 推出 $(u_1,v)=(u,v_1)$, 即
\eqref{eq:7.25}.

设 $u\in D(A^\star)$, 令 $f=u+A^\star u$.我们有
\begin{equation*}
(f,v)=(u,v+Av)\quad\forall v\in D(A),
\end{equation*}
即
\begin{equation*}
(f,J_1w)=(u,w)\quad\forall w\in H.
\end{equation*}
因此 $u=J_1f$, 从而 $u\in D(A)$.这就证明了 $D(A^\star)=D(A)$, 因此 $A$
自伴.
\end{proof}

\begin{remark}
\label{rem:7.8}
必须注意: 若 $A$ 是单调算子 (甚至是对称单调算子), 则 $A^\star$ 不必单调.
然而可以证明下述性质等价:
\begin{equation*}
A\ \text{极大单调}\iff A^\star\ \text{极大单调}
\iff A\ \text{闭},\ D(A)\ \text{稠密},\ A\ \text{与}\ A^\star\
\text{单调}.
\end{equation*}
这个结果的一个更一般的版本见问题 16.
\end{remark}

\begin{keytheorem}
\label{thm:7.7}
设 $A$ 为自伴的极大单调算子.则对每个 $u_0\in H$, 存在唯一的函数\footnote{
让我们强调定理  \ref{thm:7.4} 与定理  \ref{thm:7.7} 的区别.这里
$u_0\in H$ (而不是 $u_0\in D(A)$); 结论是 \eqref{eq:7.6} 有解, 且该解在
远离 $t=0$ 处光滑.然而, $\left|\dfrac{du}{dt}(t)\right|$ 当 $t\to0$ 时可能
「爆炸」.}
\begin{equation*}
u\in C([0,+\infty);H)\cap C^1((0,+\infty);H)
\cap C((0,+\infty);D(A))
\end{equation*}
使得
\begin{equation*}
\begin{cases}
\dfrac{du}{dt}+Au=0 & \text{在}\ (0,+\infty)\ \text{上},\\
u(0)=u_0.
\end{cases}
\end{equation*}
此外, 我们有
\begin{equation*}
\abs{u(t)}\le\abs{u_0}
\quad\text{且}\quad
\left|\frac{du}{dt}(t)\right|=\abs{Au(t)}\le\frac1t\abs{u_0}
\quad\forall t>0,
\end{equation*}
\begin{equation}
u\in C^k((0,+\infty);D(A^\ell))\quad\forall k,\ell\ \text{为整数}.
\label{eq:7.26}
\end{equation}
\end{keytheorem}

\begin{proof}
\emph{唯一性.} 设 $u$ 与 $u'$ 是两个解.由 $A$ 的单调性, 可见
$\phi(t)=\abs{u(t)-u'(t)}^2$ 在 $(0,+\infty)$ 上不增.另一方面, $\phi$ 在
$[0,+\infty)$ 上连续且 $\phi(0)=0$.因此 $\phi\equiv0$.

\emph{存在性.} 证明分为两步:

\emph{第 1 步.} 先设 $u_0\in D(A^2)$, 并设 $u$ 是定理  \ref{thm:7.4} 给出的
\eqref{eq:7.6} 的解.我们断言
\begin{equation}
\left|\frac{du}{dt}(t)\right|\le\frac1t\abs{u_0}\quad\forall t>0.
\label{eq:7.27}
\end{equation}
如命题  \ref{prop:7.6} 的证明中那样, 我们有
\begin{equation*}
J_\lambda^\star=J_\lambda\quad\text{与}\quad A_\lambda^\star=A_\lambda
\quad\forall\lambda>0.
\end{equation*}
我们回到定理  \ref{thm:7.4} 的证明中引入的近似问题:
\begin{equation}
\frac{du_\lambda}{dt}+A_\lambda u_\lambda=0\ \text{在}\ [0,+\infty)\ \text{上},
\qquad u_\lambda(0)=u_0.
\label{eq:7.28}
\end{equation}
把 \eqref{eq:7.28} 与 $u_\lambda$ 作标量积并在 $[0,T]$ 上积分, 得到
\begin{equation}
\frac12\abs{u_\lambda(T)}^2+\int_0^T(A_\lambda u_\lambda,u_\lambda)\,dt
=\frac12\abs{u_0}^2.
\label{eq:7.29}
\end{equation}
把 \eqref{eq:7.28} 与 $t\dfrac{du_\lambda}{dt}$ 作标量积并在 $[0,T]$ 上
积分, 得到
\begin{equation}
\int_0^T\left|\frac{du_\lambda}{dt}(t)\right|^2t\,dt
+\int_0^T\left(A_\lambda u_\lambda(t),
\frac{du_\lambda}{dt}(t)\right)t\,dt=0.
\label{eq:7.30}
\end{equation}
但
\begin{equation*}
\frac{d}{dt}(A_\lambda u_\lambda,u_\lambda)
=\left(A_\lambda\frac{du_\lambda}{dt},u_\lambda\right)
+\left(A_\lambda u_\lambda,\frac{du_\lambda}{dt}\right)
=2\left(A_\lambda u_\lambda,\frac{du_\lambda}{dt}\right),
\end{equation*}
因为 $A_\lambda^\star=A_\lambda$.对 \eqref{eq:7.30} 中的第二个积分作分部
积分, 得到
\begin{align}
\int_0^T\left(A_\lambda u_\lambda(t),
\frac{du_\lambda}{dt}(t)\right)t\,dt
&=\frac12\int_0^T\frac{d}{dt}\bigl[(A_\lambda u_\lambda,u_\lambda)\bigr]t\,dt
\notag\\
&=\frac12(A_\lambda u_\lambda(T),u_\lambda(T))T
-\frac12\int_0^T(A_\lambda u_\lambda,u_\lambda)\,dt.
\label{eq:7.31}
\end{align}
另一方面, 由于函数 $t\mapsto\left|\dfrac{du_\lambda}{dt}(t)\right|$ 不增
(由引理  \ref{lem:7.1}), 我们有
\begin{equation}
\int_0^T\left|\frac{du_\lambda}{dt}(t)\right|^2t\,dt
\ge\left|\frac{du_\lambda}{dt}(T)\right|^2\frac{T^2}{2}.
\label{eq:7.32}
\end{equation}
合并 \eqref{eq:7.29}, \eqref{eq:7.30}, \eqref{eq:7.31} 与
\eqref{eq:7.32}, 我们得到
\begin{equation*}
\frac12\abs{u_\lambda(T)}^2+T(A_\lambda u_\lambda(T),u_\lambda(T))
+T^2\left|\frac{du_\lambda}{dt}(T)\right|^2\le\frac12\abs{u_0}^2;
\end{equation*}
特别地推出
\begin{equation}
\left|\frac{du_\lambda}{dt}(T)\right|\le\frac1T\abs{u_0}
\quad\forall T>0.
\label{eq:7.33}
\end{equation}
最后, 在 \eqref{eq:7.33} 中令 $\lambda\to0$ 取极限.这就完成了
\eqref{eq:7.27} 的证明, 因为
$\dfrac{du_\lambda}{dt}\to\dfrac{du}{dt}$ (见定理  \ref{thm:7.4} 证明中的
第 5 步).

\emph{第 2 步.} 现在假设 $u_0\in H$.设 $(u_{0n})$ 为 $D(A^2)$ 中的序列,
使得 $u_{0n}\to u_0$ (回忆 $D(A^2)$ 在 $D(A)$ 中稠密, 而 $D(A)$ 在 $H$ 中
稠密; 因此 $D(A^2)$ 在 $H$ 中稠密).设 $u_n$ 为
\begin{equation*}
\begin{cases}
\dfrac{du_n}{dt}+Au_n=0 & \text{在}\ [0,+\infty)\ \text{上},\\
u_n(0)=u_{0n}
\end{cases}
\end{equation*}
的解.由定理  \ref{thm:7.4}, 我们知道
\begin{equation*}
\abs{u_n(t)-u_m(t)}\le\abs{u_{0n}-u_{0m}}
\quad\forall m,n,\ \forall t\ge0,
\end{equation*}
而由第 1 步
\begin{equation*}
\left|\frac{du_n}{dt}(t)-\frac{du_m}{dt}(t)\right|
\le\frac1t\abs{u_{0n}-u_{0m}}\quad\forall m,n,\ \forall t>0.
\end{equation*}
由此推出 $u_n$ 在 $[0,+\infty)$ 上一致收敛于某个极限 $u(t)$, 且
$\dfrac{du_n}{dt}(t)$ 在每个区间 $[\delta,+\infty)$ ($\delta>0$) 上一致收敛
于 $\dfrac{du}{dt}(t)$.极限函数 $u$ 满足
\begin{equation*}
u\in C([0,+\infty);H)\cap C^1((0,+\infty);H),
\qquad u(t)\in D(A)\ \forall t>0
\end{equation*}
且
\begin{equation*}
\frac{du}{dt}(t)+Au(t)=0\quad\forall t>0
\end{equation*}
(这里用到 $A$ 闭这一事实).

现在我们转向 \eqref{eq:7.26} 的证明.我们将对 $k\ge2$ 归纳证明
\begin{equation}
u\in C^{k-j}((0,+\infty);D(A^j))\quad\forall j=0,1,\dots,k.
\label{eq:7.34}
\end{equation}
假设 \eqref{eq:7.34} 对直到 $k-1$ 的阶数都成立.特别地, 我们有
\begin{equation}
u\in C((0,+\infty);D(A^{k-1})).
\label{eq:7.35}
\end{equation}
为证明 \eqref{eq:7.34}, 只须 (由定理  \ref{thm:7.5}) 验证
\begin{equation}
u\in C((0,+\infty);D(A^k)).
\label{eq:7.36}
\end{equation}
考虑 Hilbert 空间 $\tilde H=D(A^{k-1})$ 以及由
\begin{equation*}
\begin{cases}
D(\tilde A)=D(A^k),\\
\tilde A=A
\end{cases}
\end{equation*}
定义的算子 $\tilde A:D(\tilde A)\subset\tilde H\to\tilde H$.容易看出
$\tilde A$ 在 $\tilde H$ 中极大单调且对称; 因此它自伴.把定理  \ref{thm:7.7}
的第一个断言应用于空间 $\tilde H$ 中的算子 $\tilde A$, 可见对任何
$v_0\in\tilde H$, 问题
\begin{equation}
\begin{cases}
\dfrac{dv}{dt}+Av=0 & \text{在}\ (0,+\infty)\ \text{上},\\
v(0)=v_0
\end{cases}
\label{eq:7.37}
\end{equation}
有唯一解 $v$.此外,
\begin{equation*}
v\in C([0,+\infty);\tilde H)\cap C^1((0,+\infty);\tilde H)
\cap C((0,+\infty);D(\tilde A)).
\end{equation*}
取 $v_0=u(\varepsilon)$ ($\varepsilon>0$)——由 \eqref{eq:7.35} 我们已经知道
$v_0\in\tilde H$——我们推出 $u\in C((\varepsilon,+\infty);D(A^k))$, 这就
完成了 \eqref{eq:7.36} 的证明.
\end{proof}

\section*{第 7 章评注}
\addcontentsline{toc}{section}{第 7 章评注}
\markright{第 7 章评注}

\cmtitem{1. Banach 空间中的 Hille--Yosida 定理.}
Hille--Yosida 定理可以推广到 Banach 空间.确切的陈述如下.设 $E$ 为 Banach
空间, $A:D(A)\subset E\to E$ 为无界线性算子.若 $\overline{D(A)}=E$ 且对
每个 $\lambda>0$, $I+\lambda A$ 是由 $D(A)$ 到 $E$ 的双射, 并且
$\norm{(I+\lambda A)^{-1}}_{\LL(E)}\le1$, 则称 $A$ 是 \emph{$m$-增生}的.

\begin{theorem}[Hille--Yosida]
\label{thm:7.8}
设 $A$ 为 $m$-增生算子.则给定任何 $u_0\in D(A)$, 存在唯一的函数
\begin{equation*}
u\in C^1([0,+\infty);E)\cap C([0,+\infty);D(A))
\end{equation*}
使得
\begin{equation}
\begin{cases}
\dfrac{du}{dt}+Au=0 & \text{在}\ [0,+\infty)\ \text{上},\\
u(0)=u_0.
\end{cases}
\label{eq:7.38}
\end{equation}
此外,
\begin{equation*}
\norm{u(t)}\le\norm{u_0}
\quad\text{且}\quad
\left\|\frac{du}{dt}(t)\right\|=\norm{Au(t)}\le\norm{Au_0}
\quad\forall t\ge0.
\end{equation*}
把映射 $u_0\mapsto u(t)$ 按连续性延拓到整个 $E$, 记作 $S_A(t)$.它是 $E$
上的连续压缩半群.反之, 给定任何连续压缩半群 $S(t)$, 存在唯一的
$m$-增生算子 $A$, 使得 $S(t)=S_A(t)\ \forall t\ge0$.
\end{theorem}

关于证明, 见例如 P.~Lax \bref{Lax}{1}, A.~Pazy \bref{Pazy}{1}, J.~Goldstein \bref{Goldstein}{1}, E.~Davies \bref{Davies}{1}, \bref{Davies}{2}, K.~Yosida \bref{Yosida}{1}, M.~Reed--B.~Simon
\bref{ReedSimon}{1} 第 2 卷, H.~Tanabe \bref{Tanabe}{1}, N.~Dunford--J.~T.~Schwartz \bref{DunfordSchwartz}{1} 第 1 卷, M.~Schechter \bref{Schechter}{1}, A.~Friedman \bref{Friedman}{2}, R.~Dautray--J.-L.~Lions \bref{DautrayLions}{1} 第 XVII 章, A.~Balakrishnan \bref{Balakrishnan}{1}, T.~Kato \bref{Kato}{1}, W.~Rudin \bref{Rudin}{1}.这些文献对半群理论作了大量展开.

\cmtitem{2. 指数公式.}
求解 \eqref{eq:7.38} 有许多迭代技巧.这里提一种基本方法.

\begin{theorem}
\label{thm:7.9}
假设 $A$ 为 $m$-增生算子.则对每个 $u_0\in D(A)$, \eqref{eq:7.38} 的解
$u$ 由「指数公式」给出:
\begin{equation}
u(t)=\lim_{n\to+\infty}
\left[\left(I+\frac tnA\right)^{-1}\right]^nu_0.
\label{eq:7.39}
\end{equation}
\end{theorem}

证明见例如 K.~Yosida \bref{Yosida}{1} 与 A.~Pazy \bref{Pazy}{1}.公式
\eqref{eq:7.39} 用数值分析的语言来说, 对应于 \eqref{eq:7.38} 的隐式时间
离散格式的收敛性 (见例如 K.~W.~Morton--D.~F.~Mayers
\bref{MortonMayers}{1}).更确切地说, 把区间 $[0,t]$ 分成 $n$ 个等长的小区间
$\Delta t=t/n$, 并从 $u^0$ 出发递推求解方程
\begin{equation*}
\frac{u^{j+1}-u^j}{\Delta t}+Au^{j+1}=0,
\quad j=0,1,\dots,n-1.
\end{equation*}
换言之, $u^n$ 由
\begin{equation*}
u^n=(I+\Delta tA)^{-n}u_0
=\left(I+\frac tnA\right)^{-n}u_0
\end{equation*}
给出.当 $n\to\infty$ (即 $\Delta t\to0$) 时, $u^n$ 收敛于 $u(t)$ 是
「符合直觉的」.

\cmtnum{3.} 定理  \ref{thm:7.7} 是解析半群理论的第一步.关于这个专题, 见
例如 K.~Yosida \bref{Yosida}{1}, T.~Kato \bref{Kato}{1}, M.~Reed--B.~Simon \bref{ReedSimon}{1} 第 2 卷, A.~Friedman
\bref{Friedman}{2}, A.~Pazy \bref{Pazy}{1}, 以及 H.~Tanabe
\bref{Tanabe}{1}.

\cmtitem{4. 非齐次方程. 非线性方程.}
在 Banach 空间 $E$ 中考虑问题
\begin{equation}
\begin{cases}
\dfrac{du}{dt}(t)+Au(t)=f(t) & \text{在}\ [0,T]\ \text{上},\\
u(0)=u_0.
\end{cases}
\label{eq:7.40}
\end{equation}
下述结论成立.

\begin{theorem}
\label{thm:7.10}
假设 $A$ 为 $m$-增生算子.则对每个 $u_0\in D(A)$ 与每个
$f\in C^1([0,T];E)$, 存在 \eqref{eq:7.40} 的唯一解 $u$, 满足
\begin{equation*}
u\in C^1([0,T];E)\cap C([0,T];D(A)).
\end{equation*}
此外, $u$ 由公式
\begin{equation}
u(t)=S_A(t)u_0+\int_0^tS_A(t-s)f(s)\,ds
\label{eq:7.41}
\end{equation}
给出, 其中 $S_A(t)$ 是评注 1 中引入的半群.
\end{theorem}

注意, 若只假设 $f\in L^1((0,T);E)$, 公式 \eqref{eq:7.41} 仍有意义, 并
给出 \eqref{eq:7.40} 的广义解.关于这些问题, 见例如 T.~Kato
\bref{Kato}{1}, A.~Pazy \bref{Pazy}{1}, R.~H.~Martin \bref{Martin}{1}, H.~Tanabe \bref{Tanabe}{1}.

在物理应用中会遇到许多形如
\begin{equation*}
\frac{du}{dt}+Au=F(u)
\end{equation*}
的「半线性」方程, 其中 $F$ 是由 $E$ 到 $E$ 的非线性映射.关于这些问题,
见例如 R.~H.~Martin \bref{Martin}{1}, Th.~Cazenave--A.~Haraux \bref{CazenaveHaraux}{1}, 以及第~
\ref{chap:10}~章的评注.

我们还要提到, 第~
\ref{chap:7}~章的一些结果有非线性推广.考虑非线性的 $m$-增生算子
$A:D(A)\subset E\to E$ 是有用的.关于这个专题, 见例如 H.~Brezis
\bref{Brezis}{1} 与 V.~Barbu \bref{Barbu}{1}.
